\documentclass[12pt,a4paper]{article}

% Packages for margins, footnotes, and line spacing
\usepackage[left=2.5cm,right=2.5cm,top=2.5cm,bottom=2.5cm]{geometry}
\usepackage{setspace}
\usepackage{fancyhdr}
\usepackage{graphicx}
\usepackage{caption}
\usepackage{amsmath}
\usepackage{parskip}
\usepackage[ngerman]{babel}
% csquotes: \enquote{...} statt \glqq...\grqq. autostyle folgt babel (-> „…").
\usepackage[autostyle=true,german=quotes]{csquotes}
% Aktiviert gerade Anführungszeichen: "Text" wird automatisch zu „Text".
% (Verschachtelt: "Außen \enquote{innen}" -> „Außen ‚innen'".)
\MakeOuterQuote{"}
\usepackage{quoting}
\usepackage{etoolbox}
\usepackage{hanging}  % For hanging indentation
\usepackage{tabto}
\usepackage{tabularx}    % For tables that automatically fit page width
\usepackage{makecell}    % For better cell formatting
\usepackage{array}       % For advanced table features
\usepackage{booktabs}    % For professional-looking tables
\usepackage{multirow}
\usepackage{multicol}
\usepackage{mathabx}
\usepackage{xcolor}
\usepackage[T1]{fontenc}
\usepackage{lipsum}      % Blindtext (\lipsum) -- vor Abgabe entfernen


\renewcommand{\arraystretch}{1.5}

% Define a new column type for automatically wrapped and centered content
\newcolumntype{C}{>{\centering\arraybackslash}X}
\newcolumntype{R}[1]{>{\centering\arraybackslash}p{#1}}

% Command for table cell with forced line spacing of 1.0
\newcommand{\tablecell}[1]{%
	\begin{tabular}[c]{@{}c@{}}
		\setstretch{1.0}#1
	\end{tabular}%
}


%CITATION
% biblatex auto-loads biblatex.cfg from the project directory (Heidelberg
% German-linguistics style). Keep biblatex.cfg next to this file.
\usepackage[backend=biber, style=authoryear]{biblatex}
\addbibresource{references.bib}


% LINE SPACING
% Word's "1.5 lines" is a true 1.5x stretch, not setspace's \onehalfspacing
% (which is ~1.25). Use \setstretch{1.5} to match Word.
\setstretch{1.5}
\sloppy


% PARAGRAPH SETTINGS
\setlength{\parskip}{0pt}
\setlength{\parindent}{2em}


% FOOTNOTE SETTINGS
\renewcommand{\footnoterule}{%
	\kern -3pt
	\hrule width 2in height 0.4pt
	\kern 2.6pt
}
\renewcommand{\footnotesize}{\small}
\setlength{\footnotesep}{0.0cm}


% NUMMERIERTE BEISPIELE
% Counter for numbered linguistic examples (Beispiele)
\newcounter{bsp}

% Quoting setup: each quoting environment is prefixed with (1), (2), ...
\quotingsetup{
	leftmargin=2cm, rightmargin=1cm, vskip=12pt, font=small,
	begintext={(\thebsp)\tabto*{0.0cm}},
}
% Step the example counter at the start of every quoting environment
\AtBeginEnvironment{quoting}{%
	\setlength{\parindent}{-25pt}%
	\stepcounter{bsp}
}
% Reset example numbering at the start of the document
\AtBeginDocument{\setcounter{bsp}{0}}

% \bspback{n} -> number of the example n positions before the current one
% (\bspback{} = current example, i.e. n defaults to 0)
\newcommand{\bspback}[1]{\the\numexpr\value{bsp}-0#1\relax}

% SECTION NUMBER SETTINGS
\renewcommand{\thesection}{\arabic{section}}
\renewcommand{\thesubsection}{\thesection.\arabic{subsection}}
\renewcommand{\thesubsubsection}{\thesubsection.\arabic{subsubsection}}
\renewcommand{\thefigure}{\thesection.\arabic{figure}}
\renewcommand{\thetable}{\thesection.\arabic{table}}


% ===== Angaben für die Titelseite =====
\newcommand{\universitaet}{Universität Heidelberg}
\newcommand{\seminar}{Germanistisches Seminar}
\newcommand{\semester}{Wintersemester JJJJ/JJJJ}
\newcommand{\lehrveranstaltung}{Titel der Lehrveranstaltung}
\newcommand{\dozent}{Prof. Dr. Name Nachname}
\newcommand{\arbeitstitel}{Titel der Arbeit}
\newcommand{\untertitel}{Untertitel der Arbeit}
\newcommand{\autorname}{Vorname Nachname}
\newcommand{\matrikelnr}{0000000}
\newcommand{\studiengang}{Studiengang}
\newcommand{\email}{vorname.nachname@stud.uni-heidelberg.de}
\newcommand{\abgabedatum}{TT.MM.JJJJ}
% ======================================

\begin{document}

% Title page
\begin{titlepage}
	\noindent
	\universitaet \\
	\seminar \\
	\semester \\
	Lehrveranstaltung: \lehrveranstaltung \\
	Dozent: \dozent

	\centering
	\vspace*{2cm}
	\Huge
	\textbf{\arbeitstitel}

	\vspace{1cm}
	\LARGE
	\untertitel

	\vfill
	\large
	\autorname

	\vspace{0.1cm}
	\large
	Matrikelnr. \matrikelnr

	\vspace{0.1cm}
	\studiengang

	\vspace{0.1cm}
	\email

	\vspace{0.1cm}
	\abgabedatum

	\vspace*{1cm}
\end{titlepage}


% Table of contents (without page nbr.)
\thispagestyle{empty}
\tableofcontents
\newpage

% Anführungszeichen (Deutsch „…"):
%   \enquote{Text}   ODER einfach   "Text"   (statt \glqq...\grqq)
% Zitierbefehle:
%   \autocite{key}        -> (Nachname Jahr)
%   \autocite[123]{key}   -> (Nachname Jahr: 123)
%   \textcite{key}        -> Nachname (Jahr)
% Nummerierte Beispiele: die quoting-Umgebung wird automatisch fortlaufend
% nummeriert -- (1), (2), ... \bspback{n} gibt die Nummer des n-ten Beispiels
% vor dem aktuellen Zähler zurück (für Rückverweise im Fließtext).

\setcounter{page}{1}

\section{Einleitung}
% Beginn der Seitenzählung! Thematische Einführung, Problemstellung,
% Untersuchungsfrage und Ziel der Arbeit, zentrale These(n), Abgrenzung von
% benachbarten Fragestellungen.
\lipsum[1-2]


\section{Theoretischer Rahmen}
% Theoretische Konzepte aus der Forschungsliteratur, auf die zurückgegriffen
% wird; kritische Diskussion verwandter Konzepte; Einführung der zentralen
% Begriffe und der verwendeten Terminologie.
\lipsum[3-4]


\section{Korpus und Methode}
% Zentraler Teil empirischer Hausarbeiten: Erläuterung der Datengrundlage,
% fremdes oder eigenes Korpus (im letzteren Fall Vorgehen beim Datenerheben);
% Darstellung und Begründung der ausgewählten Methode.
\lipsum[5-6]


\section{Hauptteil}
% Präsentiert die Ergebnisse der Untersuchung; in mehrere Teilkapitel
% untergliedert, die jeweils aussagekräftige Überschriften enthalten und dem
% Aufbau der Argumentation folgen.

\subsection{Unterkapitel}
\lipsum[7]

% Nummerierte Beispiele aus dem Korpus:
\begin{quoting}
Dies ist ein nummeriertes Beispiel aus dem Korpus.
\end{quoting}

\begin{quoting}
Ein zweites Beispiel.
\end{quoting}

Auf das letztgenannte Beispiel~(\bspback{}) kann mit \texttt{\textbackslash bspback\{\}} verwiesen werden, auf vorige Beispiele mit \texttt{\textbackslash bspback\{n-zurückliegend\}}, z.\,B. auf Beispiel~(\bspback{1}).

\subsubsection{Unterunterkapitel}
\lipsum[8]

\subsubsection{Unterunterkapitel}
\lipsum[9]

\subsection{Unterkapitel}
\lipsum[10]

\subsubsection{Unterunterkapitel}
\lipsum[11]


\section{Fazit und Ausblick}
% Zusammenfassung der Ergebnisse mit Rückbezug auf die Leitfrage, die
% Ausgangshypothese(n), die Problemstellung in pointierter Form, ggf. mit
% Hinweisen auf weiterführende Fragen.
\lipsum[12]


\section{Bibliographie}
\printbibliography[heading=none]

\end{document}
